\chapter{Introduction}

The internship was carried out in a company engineering context where the goal
was not only to find links on suspected illegal streaming websites. For
rights-enforcement review, a useful result must show the page path, visible
context, screenshots, stream or embedded-player indicators, provider context,
and the trace of the browser actions that produced the result. The project
therefore addresses evidence collection, not simple scraping.

Premium sports broadcasting depends on exclusive distribution rights. Rights
holders and broadcasters pay to transmit competitions in specific territories,
on specific platforms, and under controlled commercial conditions. Illegal
streaming websites weaken that model by redistributing protected events without
authorization, often through domains, mirrors, and player pages that change
frequently. In that environment, technical evidence must be collected quickly
but also carefully enough to remain reviewable after the page has changed.

The first practical challenge is therefore not only discovering suspicious
websites. It is preserving enough evidence to explain what was observed:
which URL was opened, what visible page state appeared, which player or iframe
was reached, which stream or manifest indicator was detected, which provider or
channel clues were available, and which browser actions produced those results.
Screenshots, DOM structure, frame paths, network requests, stream references,
and runtime traces all become part of the evidence picture. They are often
hidden behind JavaScript interfaces, popups, nested frames, language variation,
server tabs, redirects, and short-lived media URLs.

This turns the internship problem into an engineering problem rather than a
simple scraping task. Manual review is slow and inconsistent when the number of
targets grows. Static extraction is too fragile for modern streaming pages
because the useful evidence may appear only after the page loads scripts,
reacts to clicks, switches servers, or redirects through several layers. The
report therefore starts from the business need, then progressively narrows
toward the browser-agent artifact designed to make this evidence collection
more traceable and maintainable.

\section{Host Organization and Internship Context}

\subsection{Host Organization}

Soft Stars is a software company working in digital media intelligence and
technical support for online content-protection operations. Its work focuses
on helping media organizations observe online distribution channels, identify
unauthorized use of protected content, and prepare technical material that can
be reviewed by operational, legal, or business teams. The PFE work was
connected to the beIN SPORTS rights-enforcement context, where reviewers need
technical evidence that can be inspected before any operational or legal action
is considered.

\subsection{Existing Technical Context}

Before this internship, the company already relied on Puppeteer-based crawlers
to inspect streaming websites. These crawlers were useful on known website
patterns: they could visit prepared targets, follow expected selectors, collect
candidate links, and support established evidence-collection flows. Their main
limit was maintenance. Small website changes could break selectors, navigation
paths, popup behavior, iframe discovery, server switching, or stream-detection
logic. The internship therefore explored a more traceable and adaptable
browser-agent workflow that could support the existing crawler environment.

\subsection{Operational Evidence-Collection Context}

Sports-streaming investigations are difficult because evidence is often hidden
behind dynamic pages, scripts, popups, server buttons, redirects, and nested
iframes. A suspicious website can expose one page for discovery, another page
for hosting the visible player, and a separate embedded player URL that finally
produces media requests. The evidence path therefore matters as much as the
final stream candidate.

For a rights-enforcement team, a final URL by itself is not enough. The
reviewer needs to understand how the browser reached the page, whether the
player was visible, whether the stream indicator appeared after a valid action,
and whether screenshots or runtime traces support the result. The same website
can behave differently by date, event, country, browser profile, mirror, or
server tab. A traceable workflow makes those differences explicit instead of
hiding them behind a single success or failure label.

The system also needs to identify visible channel or broadcaster context when
the page exposes it. The workflow is connected to beIN SPORTS, but the evidence
model is not hardcoded only for beIN SPORTS. It attempts to detect channel
signals from the page, screenshot, DOM text, stream context, and metadata, then
uses that context later when preparing reviewable email drafts.

\section{Problem Statement}

\subsection{Existing Crawler Limitations}

The existing crawler baseline solved useful collection tasks on known layouts,
but it was fragile when websites changed. Selectors could disappear, player
buttons could move, server labels could change, popups could steal focus, and a
previously direct iframe could become hidden behind a new interaction. Each
change created a maintenance task before evidence could be collected reliably
again.

This limitation does not mean that the previous crawlers were useless. On
stable patterns, a Puppeteer crawler is fast, deterministic, and easier to run
at scale than a fully adaptive agent. The problem appears when the crawler's
assumptions no longer match the live site. In those cases the system needs a
way to observe the current rendered page, reason about the next browser action,
verify the result, and preserve a trace that explains the failure or recovered
evidence.

\subsection{Limits of Static or Rule-Based Extraction}

Static scraping misses JavaScript-rendered players, delayed iframes, post-click
media requests, and tokenized manifests. A simple HTTP request can miss the
real player state, while manual review becomes slow and inconsistent when the
number of targets increases. Rule-based automation also struggles when a page
requires observation, action, and verification rather than one fixed selector
chain.

\subsection{Need for Traceable Browser-Agent Evidence Collection}

The problem addressed in this report is the following: how can a browser-based
system support evidence collection from dynamic illegal streaming websites
while keeping the investigation traceable, repeatable, and measurable? The
answer explored during the internship is a browser-agent artifact where each
specialist owns a bounded part of the task and where screenshots, network
signals, model calls, tool calls, and final outputs are preserved for review.

\input{figures/tikz/fig-01-crawler-limitation}

\section{Internship Mission and Scope}

\subsection{Assigned Mission}

My mission was to study the evidence-collection workflow, analyze the limits of
the existing crawler context, validate whether an LLM-guided browser agent
could operate on real streaming pages, and design Open Web Catcher (OWC) as a
maintainable reference architecture for that approach.

\subsection{Mission Deliverables}

The internship produced three connected deliverables. The first was an analysis
of the crawler maintenance problem and the evidence expected by reviewers. The
second was an n8n landing-agent workflow used to validate browser-agent
behavior on real targets and deployed as a backup solution to support the
existing crawler workflow when crawler rules failed to recover useful
landing-page candidates. The third was OWC, a modular browser-agent
architecture with an orchestrator, classification, landing, hosting, embedded,
provider, and email-review stages.

\subsection{Scope Boundaries}

The work remained focused on evidence collection and review support. It did not
send takedown emails automatically, it did not download protected content, it
did not claim legal certainty from technical signals alone, and it did not
claim that the full Python/FastAPI/Puppeteer platform replaced the operational
company crawler. OWC should be read as a maintainable browser-agent
architecture and evaluation foundation.

\section{Objectives and Main Contributions}

\subsection{Project Objectives}

The project objectives were to classify pages, route work to specialist agents,
interact with browser state, capture screenshots, preserve manifests and media
indicators, detect channel or broadcaster context where possible, enrich
provider metadata, store traces, and expose reviewable outputs for a human
operator.

These objectives were defined from the real evidence path. A landing page must
be inspected for event cards and candidate hosting links. A hosting page must
be able to deal with play controls, server buttons, popups, redirects, and
iframes. An embedded player page must be inspected with a narrower focus on
player state and stream indicators. Provider and channel enrichment then adds
context that helps reviewers understand who may be involved and which rights
holder context may be relevant.

\subsection{Engineering Contributions}

The engineering contributions are a browser-agent architecture, a clear browser
tool layer, specialist prompts with explicit handoffs, a trace model for
successful and failed runs, and a persistence design that stores both evidence
and runtime events. These contributions make future maintenance easier because
a failure can be tied to a route decision, browser action, tool result,
provider lookup, or persistence issue.

\subsection{Review and Evaluation Contributions}

The review contribution is an evidence bundle that groups screenshots, stream
or embedded-player candidates, channel/provider context, tool calls, model
usage, and reviewable email drafts. The evaluation contribution is a
regression-oriented approach that compares browser-agent behavior across
website batches instead of relying on one successful demonstration.

\section{Report Organization}

The report follows a general-to-specific structure. Chapter~2 presents DSRM as
the sole methodology and defines requirements and constraints. Chapter~3
introduces streaming-page roles, browser-rendered evidence, media manifests,
browser automation, and traceability. Chapter~4 presents the architecture,
including the browser tool layer and agent responsibility views. Chapter~5
describes the implementation from the existing crawler baseline and n8n backup
workflow to OWC. Chapter~6 evaluates the workflow through regression results,
tool behavior, cost telemetry, provider/channel enrichment, and representative
cases. Chapter~7 gives the general conclusion.

\begin{figure}[!htbp]
\centering
\resizebox{0.96\textwidth}{!}{%
\begin{tikzpicture}[
  chapter/.style={draw=accent!75, fill=blue!5, rounded corners=6pt,
                  minimum width=3.2cm, minimum height=1.18cm,
                  text width=2.9cm, align=center,
                  font=\footnotesize\bfseries, inner sep=5pt},
  core/.style={draw=orange!75!black, fill=orange!9, rounded corners=6pt,
               minimum width=3.2cm, minimum height=1.18cm,
               text width=2.9cm, align=center,
               font=\footnotesize\bfseries, inner sep=5pt},
  eval/.style={draw=green!55!black, fill=green!8, rounded corners=6pt,
               minimum width=3.2cm, minimum height=1.18cm,
               text width=2.9cm, align=center,
               font=\footnotesize\bfseries, inner sep=5pt},
  close/.style={draw=purple!60!black, fill=purple!7, rounded corners=6pt,
                minimum width=3.2cm, minimum height=1.18cm,
                text width=2.9cm, align=center,
                font=\footnotesize\bfseries, inner sep=5pt},
  band/.style={draw=gray!35, fill=gray!3, rounded corners=9pt, inner xsep=11pt, inner ysep=8pt},
  arrow/.style={->, >=Stealth, thick, accent!72, shorten <=3pt, shorten >=3pt},
  stage/.style={font=\footnotesize\bfseries, text=gray!75, fill=white, inner sep=2pt},
  note/.style={font=\scriptsize, text=gray!70, align=center}
]

% Framing row
\node[chapter] (c1) at (0,3.25) {Ch.1\\Introduction};
\node[chapter] (c2) at (4.0,3.25) {Ch.2\\Methodology\\and Requirements};
\node[chapter] (c3) at (8.0,3.25) {Ch.3\\Technical\\Background};

% Contribution row, arranged as a reading snake to avoid crossed arrows
\node[core] (c4) at (8.0,0.85) {Ch.4\\System\\Architecture};
\node[core] (c5) at (4.0,0.85) {Ch.5\\Implementation};
\node[eval] (c6) at (0,0.85) {Ch.6\\Evaluation\\and Testing};

% Closing chapter
\node[close] (c7) at (4.0,-1.35) {Ch.7\\General\\Conclusion};

\begin{scope}[on background layer]
  \node[band, fit=(c1)(c2)(c3)] (general) {};
  \node[band, fit=(c4)(c5)(c6)] (specific) {};
\end{scope}

\node[stage, anchor=south west] at ([xshift=0.22cm,yshift=0.10cm]general.north west)
  {General framing: problem, domain, method};
\node[stage, anchor=south west] at ([xshift=0.22cm,yshift=0.10cm]specific.north west)
  {Specific contribution: architecture, implementation, validation};

\draw[arrow] (c1) -- (c2);
\draw[arrow] (c2) -- (c3);
\draw[arrow] (c3.south) -- (c4.north);
\draw[arrow] (c4) -- (c5);
\draw[arrow] (c5) -- (c6);
\draw[arrow] (c6.south) |- (c7.west);

\node[note] at ([yshift=-0.35cm]c7.south)
  {The report connects the internship problem, design choices, implementation, evaluation, and final perspectives.};
\end{tikzpicture}%
}
\caption{Report reading map connecting the internship context, design, implementation, evaluation, and conclusion.}
\label{fig:report-map}
\end{figure}

